\section*{Capítulo \ref{ch:b2:planet-formation-meteorites} --- La formación de los planetas y los meteoritos}

\begin{solution}{exo:b2:planet-formation-meteorites:1}
Gas: \qty{2200}{g.cm^{-2}} a \qty{1}{au}; a \qty{5.2}{au},
$2200\times5.2^{-3/2} = 2200/11.86 = \qty{186}{g.cm^{-2}}$. Sólidos:
$33\times0.33 = \qty{10.9}{g.cm^{-2}}$ a \qty{1}{au};
$33/11.86 = \qty{2.78}{g.cm^{-2}}$ a \qty{5.2}{au}. A \qty{5.2}{au} el gas es
$186/2.78 = 67$ veces más pesado que los sólidos.
\end{solution}

\begin{solution}{exo:b2:planet-formation-meteorites:2}
Aluminio (\qty{1653}{K}, ultrarrefractario), magnesio (\qty{1336}{K}, común),
hierro (\qty{1334}{K}, común), sodio (\qty{958}{K}, moderadamente volátil), azufre
(\qty{664}{K}, volátil), oxígeno como hielo de agua (\qty{180}{K}, muy volátil).
Las CAI, los primeros condensados, están enriquecidas en aluminio (y en calcio y titanio);
la Tierra está empobrecida en los volátiles, el sodio, el azufre y el agua, respecto al
magnesio y al aluminio.
\end{solution}

\begin{solution}{exo:b2:planet-formation-meteorites:3}
Caídas: \omterm{def:b2:planet-formation-meteorites:chondrule}{condritas} ordinarias $856/1107 = \qty{77.3}{\percent}$, \omterm{def:b2:planet-formation-meteorites:classes}{meteoritos metálicos}
$49/1107 = \qty{4.4}{\percent}$. Hallazgos: $38647/44534 = \qty{86.8}{\percent}$ y
$1020/44534 = \qty{2.3}{\percent}$. Hallazgos antárticos: $22096/44534 =
\qty{49.6}{\percent}$. Las caídas: toda caída observada se recoge sea cual sea su
clase, mientras que un hallazgo depende de lo bien que el \omterm{def:b2:planet-formation-meteorites:fall-find}{meteorito} sobrevive en el suelo,
de su aspecto y de dónde se busca, y la mitad de los hallazgos proceden de un solo tipo de
lugar.
\end{solution}

\begin{solution}{exo:b2:planet-formation-meteorites:4}
$\Delta t = 1.035\ln(\num{5.23e-5}/\num{1.0e-5}) = 1.035\ln 5.23 =
1.035\times1.654 = \qty{1.71}{Myr}$ después de las CAI.
\end{solution}

\begin{solution}{exo:b2:planet-formation-meteorites:5}
$4\pi\Sigma_1 r_1^2 = 4\pi\times\num{2.2e4}\times(\num{1.496e11})^2 =
\qty{6.19e27}{kg}$. Dentro de \qty{30}{au}: $\times\sqrt{30} = 5.48$, es decir
\qty{3.39e28}{kg} $= 0.017$ masas solares; dentro de \qty{100}{au}: $\times10$,
\qty{6.19e28}{kg} $= 0.031$ masas solares. Ambas están dentro del intervalo observado,
más cerca de su extremo superior: la nebulosa era un disco corriente.
\end{solution}

\begin{solution}{exo:b2:planet-formation-meteorites:6}
$r = (120/T)^{7/3}$ au: para el agua, $(120/182)^{7/3} = 0.659^{7/3} =
\qty{0.38}{au}$; para el metano, $(120/41)^{7/3} = 2.93^{7/3} = \qty{12.3}{au}$.
El grano bajo luz solar directa es la capa superficial del disco; el plano medio,
protegido por el polvo que tiene encima, es más frío, de modo que el hielo sobrevive más cerca. Los
\omterm{def:b2:planet-formation-meteorites:planetesimal}{planetesimales} se formaron en el plano medio, así que se aplica la estimación del plano medio, siempre que
el disco sea pasivo; el calentamiento por acreción desplaza la \omterm{def:b2:planet-formation-meteorites:snow-line}{línea de nieve} hacia fuera.
\end{solution}

\begin{solution}{exo:b2:planet-formation-meteorites:7}
A \qty{1}{au}: $r^2 = \num{2.238e22}$, $\Sigma_p = \qty{109}{kg.m^{-2}}$, de modo que
$2\pi b r^2\Sigma_p = 31.42\times\num{2.238e22}\times109 = \num{7.66e25}$;
$(\num{7.66e25})^{3/2} = \num{6.71e38}$, dividido por
$(3\times\num{1.989e30})^{1/2} = \num{2.443e15}$: \qty{2.75e23}{kg}
$= 0.046$ masas terrestres. A \qty{5.2}{au}: $r^2 = \num{6.05e23}$,
$\Sigma_p = \qty{27.8}{kg.m^{-2}}$, $2\pi b r^2\Sigma_p = \num{5.29e26}$,
$M = \num{1.22e40}/\num{2.443e15} = \qty{4.98e24}{kg} = 0.83$ masas terrestres.
Con \num{3.5} veces más sólidos, $M$ crece como $3.5^{3/2} = 6.55$: 5.5 masas
terrestres.
\end{solution}

\begin{solution}{exo:b2:planet-formation-meteorites:8}
Por kilogramo, $\num{1.01e19}$ átomos, de modo que
$\num{1.01e19}\times\num{5.00e-13} = \qty{5.05e6}{J.kg^{-1}}$. El aumento de
temperatura sería $\num{5.05e6}/837 = \qty{6030}{K}$. Calentar de \num{250} a
\qty{1700}{K} requiere $837\times1450 = \qty{1.21e6}{J.kg^{-1}}$: el aluminio-26
suministraba unas cuatro veces eso.
\end{solution}

\begin{solution}{exo:b2:planet-formation-meteorites:9}
$\num{5.07e10}/\num{7.49e7} = 677$. La vida media del hierro-60 es
$2.62/\ln2 = \qty{3.78}{Myr}$. Potencias iniciales por gramo:
$\num{5.07e10}/1.035 = \num{4.90e10}$ y $\num{7.49e7}/3.78 = \num{1.98e7}$ erg por
gramo y por millón de años, un cociente de unos \num{2470}: el aluminio-26 dominaba con mucho.
\end{solution}

\begin{solution}{exo:b2:planet-formation-meteorites:10}
Las potencias son iguales cuando $\num{4.90e10}\,e^{-t/1.035} = \num{1.98e7}\,e^{-t/3.78}$,
es decir $t(1/1.035 - 1/3.78) = \ln(\num{4.90e10}/\num{1.98e7}) = \ln 2475 = 7.81$;
con $0.966 - 0.265 = 0.702$ por millón de años, $t = \qty{11.1}{Myr}$. Para entonces al
hierro-60 le quedan $\num{7.49e7}\,e^{-11.1/3.78} = \num{7.49e7}\times0.053 =
\qty{4.0e6}{erg.g^{-1}}$, es decir \qty{400}{J.kg^{-1}}: miles de veces demasiado
poco para fundir nada.
\end{solution}

\begin{solution}{exo:b2:planet-formation-meteorites:11}
$\tau = 1.035\times\num{3.156e13} = \qty{3.267e13}{s}$ y
$H_0\tau = \qty{6.21e6}{J.kg^{-1}}$. Para \qty{1200}{K}:
$t_a^* = 1.035\ln\bigl(\num{6.21e6}/(837\times950)\bigr) = 1.035\ln7.81 =
\qty{2.13}{Myr}$; para \qty{1700}{K}: $1.035\ln\bigl(\num{6.21e6}/(837\times1450)\bigr)
= 1.035\ln5.12 = \qty{1.69}{Myr}$. Reducir $H_0$ a la mitad resta $\tau\ln2 =
\qty{0.72}{Myr}$ a ambos, ya que $t_a^* = \tau\ln H_0 + \tau\ln(\tau/c(T_c -
T_0))$ y solo cambia el primer término: el desplazamiento es un periodo de semidesintegración, sea cual sea la
temperatura crítica.
\end{solution}

\begin{solution}{exo:b2:planet-formation-meteorites:12}
$\kappa = 2.1/(3300\times837) = \qty{7.60e-7}{m^2.s^{-1}}$ y
$\sqrt{\kappa\tau} = \sqrt{\num{7.60e-7}\times\num{3.267e13}} = \qty{4980}{m}$,
unos \qty{5}{km}. El centro de un cuerpo de \qty{100}{km} de diámetro está a \qty{50}{km} de profundidad,
diez veces esa distancia: su calor no puede escapar mientras se desintegra el aluminio-26, y la
proposición se aplica. Un cuerpo de \qty{2}{km} pierde su calor tan deprisa como lo
produce. La fusión exige un radio de unas cuantas veces \qty{5}{km}, es decir un diámetro de unos
\qtyrange{10}{20}{km} o más, como encuentran los modelos.
\end{solution}

\begin{solution}{pb:b2:planet-formation-meteorites:1}
\textbf{1.} Una caída: se vio su llegada y se recogió a la mañana siguiente.
No está meteorizada, y se suma a la muestra no sesgada de lo que llega a la Tierra.

\textbf{2.} $856/1107 = 0.77$: unas tres posibilidades de cada cuatro.

\textbf{3.} \omterm{def:b2:planet-formation-meteorites:chondrule}{Cóndrulos} de \qty{0.3}{mm}, un \qty{9}{\percent} de metal y un olivino con
un \qty{18}{\percent} de fayalita: una \omterm{def:b2:planet-formation-meteorites:chondrule}{condrita} H (\qty{10}{\percent} de metal, Fa 16--20).

\textbf{4.} Los \omterm{def:b2:planet-formation-meteorites:chondrule}{cóndrulos} difuminados y los minerales uniformes indican metamorfismo térmico: tipos
5 a 6. Las \omterm{def:b2:planet-formation-meteorites:chondrule}{condritas} ordinarias alcanzaron como mucho unos \qty{1220}{K}, es decir
\qty{947}{\degreeCelsius}.

\textbf{5.} Las \omterm{def:b2:planet-formation-meteorites:chondrule}{condritas} ordinarias solo contienen un \qtyrange{0.1}{1}{\percent} de CAI en
volumen. La piedra está hecha de \omterm{def:b2:planet-formation-meteorites:chondrule}{cóndrulos} y matriz ensamblados después de las CAI, que
son los sólidos más antiguos: toda roca que las contiene es más joven.

\textbf{6.} Cada mineral recibió al formarse el mismo magnesio más aluminio-26
en proporción a su aluminio; desde entonces el aluminio-26 se ha convertido en
magnesio-26, de modo que cuanto más aluminio tenía un mineral por magnesio, más
magnesio-26 ganó.

\textbf{7.} $(0.139590 - 0.139500)/(15.0 - 0.05) = \num{0.000090}/14.95 =
\num{6.02e-6}$: la razón $\ce{^{26}Al}/\ce{^{27}Al}$ del \omterm{def:b2:planet-formation-meteorites:chondrule}{cóndrulo} cuando se
formó.

\textbf{8.} $0.139500 + \num{6.02e-6}\times3.00 = 0.139518$: en la recta.

\textbf{9.} $\Delta t = 1.035\ln(\num{5.23e-5}/\num{6.02e-6}) = 1.035\ln8.69 =
1.035\times2.162 = \qty{2.24}{Myr}$.

\textbf{10.} $4567.5 - 2.2 = \qty{4565.3}{\Ma}$.

\textbf{11.} $\num{5.23e-5}\,e^{-10/1.035} = \num{5.23e-5}\times\num{6.4e-5} =
\num{3.3e-9}$, demasiado poco magnesio-26 para medirlo. El hafnio-182 (periodo
\qty{8.9}{Myr}) o el manganeso-53 (\qty{3.7}{Myr}) convienen a acontecimientos unos diez millones de años
posteriores a las CAI.

\textbf{12.} $\kappa = 2.1/(3300\times837) = \qty{7.60e-7}{m^2.s^{-1}}$;
$R^2/\kappa = \num{4e8}/\num{7.60e-7} = \qty{5.26e14}{s} = \qty{16.7}{Myr}$, dieciséis
veces la vida media del aluminio-26: el centro conserva su calor.

\textbf{13.} $\num{1.9e-7}\,e^{-2.4/1.035} = \num{1.9e-7}\times0.0984 =
\qty{1.87e-8}{W.kg^{-1}}$.

\textbf{14.} $\num{1.87e-8}\times\num{3.267e13} = \qty{6.11e5}{J.kg^{-1}}$.

\textbf{15.} $\Delta T = \num{6.11e5}/837 = \qty{730}{K}$; máximo
$250 + 730 = \qty{980}{K}$.

\textbf{16.} Por encima de \qty{750}{K}, así que la roca se sinterizó y se metamorfizó; por debajo de
\qty{1200}{K}, así que no se fundieron ni el metal ni el sulfuro. La piedra es una \omterm{def:b2:planet-formation-meteorites:chondrule}{condrita} calentada,
de \omterm{def:b2:planet-formation-meteorites:type}{tipo petrológico} alto, por debajo de los \qty{1220}{K} de la pregunta 4: coherente.

\textbf{17.} $1.035\ln\bigl(\num{6.21e6}/(837\times950)\bigr) = \qty{2.13}{Myr}$:
unos \qty{0.3}{Myr} antes.

\textbf{18.} El aumento es $\Delta T(1 - e^{-(t - t_a)/\tau})$; el \qty{90}{\percent}
exige $e^{-(t - t_a)/\tau} = 0.1$, $t - t_a = 1.035\ln10 = \qty{2.38}{Myr}$, es decir
\qty{4.8}{Myr} después de las CAI.

\textbf{19.} $\Delta T = 7420\,e^{-0.8/1.035} = 7420\times0.462 = \qty{3430}{K}$,
máximo \qty{3680}{K}, muy por encima de los \qty{1700}{K} a los que los silicatos están fundidos.
El cuerpo se fundió (el calor de fusión contuvo su temperatura) y se
diferenció en un núcleo de metal y un manto y una corteza de silicatos.

\textbf{20.} $e^{(2.4 - 0.8)/1.035} = e^{1.546} = 4.7$ veces más.

\textbf{21.} El cinc y el plomo condensan hacia \qty{727}{K}, moderadamente volátiles; el magnesio
(\qty{1336}{K}) y el aluminio (\qty{1653}{K}) condensan mucho antes. Los sólidos de
la región de la piedra se reunieron antes de que el gas se enfriara lo bastante para que todo el cinc y el
plomo condensaran.

\textbf{22.} $33\times0.33\times2.5^{-3/2} = 10.89/3.95 = \qty{2.76}{g.cm^{-2}}$,
\qty{27.6}{kg.m^{-2}}.

\textbf{23.} $r^2 = (3.74\times10^{11})^2 = \num{1.40e23}$;
$2\pi\times5\times\num{1.40e23}\times27.6 = \num{1.21e26}$;
$(\num{1.21e26})^{3/2}/\num{2.443e15} = \qty{5.45e23}{kg}$, $0.091$ masas terrestres.
El cuerpo progenitor, $\tfrac43\pi(\num{2e4})^3\times3300 = \qty{1.1e17}{kg}$, es cinco
millones de veces más ligero: un \omterm{def:b2:planet-formation-meteorites:planetesimal}{planetesimal}, no un embrión.

\textbf{24.} El centro del cuerpo progenitor alcanzó un máximo de unos
\textbf{\qty[text-series-to-math=true]{707}{\degreeCelsius}} ($980 - 273$): lo bastante caliente para
recristalizar la \omterm{def:b2:planet-formation-meteorites:chondrule}{condrita}, demasiado frío para fundir su metal.
\end{solution}
